% TODO make hw stylesheet
\documentclass[11pt]{scrartcl}
\usepackage[a4paper, total={6in, 8in}]{geometry}
\usepackage[usenames,dvipsnames,svgnames]{xcolor}
\usepackage[shortlabels]{enumitem}
\usepackage[framemethod=TikZ]{mdframed}
\usepackage{amsmath,amssymb,amsthm}
\usepackage[colorlinks]{hyperref}
\usepackage{multicol}
\usepackage{tabularx}

\hypersetup{
colorlinks=true,
linkcolor=blue,
filecolor=magenta,      
urlcolor=magenta,
pdftitle={MAT 319 HW}
}

\usepackage{microtype}
\usepackage{mathtools}
\usepackage[headsepline]{scrlayer-scrpage}
\usepackage{thmtools}
\usepackage{todonotes}

\mdfdefinestyle{mdgreenbox}{linecolor=ForestGreen,backgroundcolor=ForestGreen!5,
linewidth=2pt,rightline=false,leftline=true,topline=false,bottomline=false,}
\declaretheoremstyle[headfont=\bfseries\sffamily\color{ForestGreen!70!black},
mdframed={style=mdgreenbox},headpunct={ --- },]{thmgreenbox}

\declaretheorem[style=thmgreenbox,name=Claim,numbered=no]{claim*}
\declaretheorem[style=thmgreenbox,name=Lemma,numbered=no]{lemma*}

\addtolength{\textheight}{3.14cm}
\providecommand{\ol}{\overline}
\providecommand{\ul}{\underline}
\providecommand{\eps}{\varepsilon}
\providecommand{\half}{\frac{1}{2}}
\providecommand{\dang}{\measuredangle} %% Directed angle
\providecommand{\CC}{\mathbb C}
\providecommand{\FF}{\mathbb F}
\providecommand{\NN}{\mathbb N}
\providecommand{\QQ}{\mathbb Q}
\providecommand{\RR}{\mathbb R}
\providecommand{\ZZ}{\mathbb Z}
\providecommand{\dg}{^\circ}
\providecommand{\ii}{\item}
\providecommand{\alert}[1]{{\sffamily\textbf{\textcolor{blue}{#1}}}}

\reversemarginpar

\theoremstyle{problem}
\newtheorem{innercustomthm}{Problem}
\newenvironment{problem}[1]
{\renewcommand\theinnercustomthm{#1}\innercustomthm}
{\endinnercustomthm}

\theoremstyle{remark}
\newtheorem{remarknumbered}{Remark}
\newenvironment{remark}
{\renewcommand\theremarknumbered{\hspace{-\the\fontdimen2\font\space}}\remarknumbered}
{\endremarknumbered}

\newenvironment{soln}{\begin{proof}[Solution]}{\end{proof}}
\newenvironment{walkthrough}{\noindent \textbf{Walkthrough.}}{\medskip}
\newlist{walk}{enumerate}{3}
\setlist[walk]{label=\bfseries (\alph*)}

\providecommand{\pow}{\operatorname{Pow}}
\providecommand{\vocab}[1]{{\textbf{\textcolor{ForestGreen}{#1}}}}
\providecommand{\lcm}{\operatorname{lcm}}
\providecommand{\abs}[1]{\left|#1\right|}

\usepackage{asymptote}
\begin{asydef}
defaultpen(fontsize(10pt));
unitsize(1cm);
size(8cm); // set a reasonable default
usepackage("amsmath");
usepackage("amssymb");
settings.tex="pdflatex";
settings.outformat="pdf";
// Replacement for olympiad+cse5 which is not standard
import geometry;
// recalibrate fill and filldraw for conics
void filldraw(picture pic = currentpicture, conic g, pen fillpen=defaultpen, pen drawpen=defaultpen)
{ filldraw(pic, (path) g, fillpen, drawpen); }
void fill(picture pic = currentpicture, conic g, pen p=defaultpen)
{ filldraw(pic, (path) g, p); }
// some geometry
pair foot(pair P, pair A, pair B) { return foot(triangle(A,B,P).VC); }
pair orthocenter(pair A, pair B, pair C) { return orthocentercenter(A,B,C); }
pair centroid(pair A, pair B, pair C) { return (A+B+C)/3; }
// cse5 abbreviations
path CP(pair P, pair A) { return circle(P, abs(A-P)); }
path CR(pair P, real r) { return circle(P, r); }
pair IP(path p, path q) { return intersectionpoints(p,q)[0]; }
pair OP(path p, path q) { return intersectionpoints(p,q)[1]; }
path Line(pair A, pair B, real a=0.6, real b=a) { return (a*(A-B)+A)--(b*(B-A)+B); }
// cse5 more useful functions
picture CC() {
	picture p=rotate(0)*currentpicture;
	currentpicture.erase();
	return p;
}
pair MP(Label s, pair A, pair B = plain.S, pen p = defaultpen) {
	Label L = s;
	L.s = "$"+s.s+"$";
	label(L, A, B, p);
	return A;
}
pair Drawing(Label s = "", pair A, pair B = plain.S, pen p = defaultpen) {
	dot(MP(s, A, B, p), p);
	return A;
}
path Drawing(path g, pen p = defaultpen, arrowbar ar = None) {
	draw(g, p, ar);
	return g;
}
\end{asydef}

\usepackage{mathrsfs}

\ihead{\footnotesize MAT 319 HW05}
\ohead{\footnotesize October 2, 2024}

\title{\vspace{-2em}MAT 319 HW05}
\author{\large Aditya Pahuja}
\date{\large Due 9/25}
\begin{document}
\maketitle

\begin{problem}{15.1}
	Determine which of the following series converge.
	\begin{enumerate}[\emph{\textbf{(\alph*)}}]
		\ii $\sum \frac{(-1)^n}{n}$
		\ii $\sum \frac{(-1)^nn!}{2^n}$
	\end{enumerate}
\end{problem}

\begin{soln}
	The first series converges. Since $a_n =\frac 1n$ is a nonnegative,
	decreasing sequence with limit zero, the alternating series theorem
	shows us that the series
	\[\sum_{n=1}^\infty (-1)^{n+1}a_n = \sum_{n=1}^\infty\frac{(-1)^{n+1}}{n}\]
	converges. We can pull out a constant term of $-1$ to see that
	\[\sum_{n=1}^\infty \frac{(-1)^n}{n}\]
	also converges, as desired.
	
	The second series does not converge. To prove this, we show that
	the sequence $b_n = \frac{(-1)^nn!}{2^n}$ does not converge to zero.
	\begin{claim*}
		For all natural numbers $n$, $|b_n| \ge \half$.
	\end{claim*}
	\begin{proof}
		We induct on $n$, with the base case $n = 1$ giving
		$|b_n| = |-\half| \ge \half$.
		
		In general, if the claim holds for $n$, then $n + 1 \ge 2$, so
		the chain of inequalities
		\[|b_{n+1}| = \abs{\frac{(-1)^{n+1}(n+1)!}{2^{n+1}}}
		= \abs{\frac{-1\cdot(n+1)}{2}\cdot \frac{(-1)^nn!}{2^n}}
		= \abs{\frac{-(n+1)}{2}}\cdot \abs{b_n} \ge \frac 22\cdot\abs {b_n}
		= \abs{b_n}\]
		implies that the claim then holds for $n+1$, too.
	\end{proof}
	
	In order for $\lim b_n = 0$ to hold, we need
	$\abs{b_n - 0} = \abs b_n < \eps$ eventually
	for any given positive real $\eps$; however, the claim above shows that
	the above statement fails when $\eps = \half$.
	However, $\sum b_n$ converging would imply $\lim b_n = 0$,
	so $\sum b_n$ cannot converge.
\end{soln}

\pagebreak
\begin{problem}{15.3}
	Show that $\sum_{n=2}^\infty \frac{1}{n(\log n)^p}$ converges
	if and only if $p > 1$.
\end{problem}

\begin{proof}[Solution 1]
	Since we haven't properly defined integration yet
	(or logs, but I can't avoid those), I decided to find a solution
	without integration.
	
	By factoring out the constant $\frac1{(\log 2)^p}$,
	we see that it is equivalent to determine when
	\[\sum_{n = 2}^\infty \frac{1}{n(\frac{\log n}{\log 2})^p}
	= \sum_{n = 2}^\infty \frac{1}{n(\log_2 n)^p}\]
	converges.
	Let $a_n = \frac{1}{n(\log_2 n)^p}$.
	
	First, suppose $p > 1$. For $n\ge 2$, define the sequence
	\[b_n = \frac{1}{2^r\cdot r^p}\]
	where $r$ is the integer such that $2^{r} \le n < 2^{r+1}$.
	Clearly $r\le \log_2 n$, so we can say that
	\[a_n = \frac 1n \cdot \frac{1}{(\log_2 n)^p} \le
	\frac{1}{2^r\cdot r^p} = b_n.\]
	Since $a_n$ and $b_n$ are both nonnegative
	and $a_n \le b_n$ for all $n\ge 2$, the comparison test says that
	the convergence of $\sum_{n\ge 2} a_n$ follows from
	the convergence of $\sum_{n\ge 2} b_n$,
	so we now prove the convergence of $\sum_{n\ge 2}b_n$.
	
	Since $b_n > 0$ for all $n$, the sequence $(s_n)$ of
	partial sums of $(b_n)$ has a limit $L$,
	which may be either a positive real number or $+\infty$.
	Moreover, any subsequence of $(s_n)$ converges to $L$,
	so consider the subsequence $(s_{n_k})_{k\ge 2}$ given by
	$n_k = 2^k - 1$ for integers $k\ge 2$.
	We can compute
	\begin{align*}
		s_{2^k - 1} &= \sum_{j=2}^{2^k-1} b_j
		= \sum_{r=1}^{k-1} \sum_{j=2^r}^{2^{r+1}-1} b_j\\
		&= \sum_{r=1}^{k-1} \sum_{j=2^r}^{2^{r+1}-1} \frac{1}{2^r\cdot r^p}
		= \sum_{r=1}^{k-1} \frac{1}{2^r\cdot r^p}\cdot 2^r\\
		&= \sum_{r=1}^{k-1} \frac{1}{ r^p}.
	\end{align*}
	The first line comes from partitioning the set
	$\{2, 3, \dots, 2^k-1\}$ into subsets of the form
	$\{2^r, 2^r + 1, \dots, 2^{r+1} - 1\}$; that is,
	the set corresponding to $r$
	is the set of integers $x$ satisfying $2^r\le x < 2^{r+1}$,
	and each $x\in\{2, 3, \dots, 2^{k} - 1\}$ satisfies
	$2^r \le x < 2^{r+1}$ for exactly one integer $r$
	between $1$ and $k - 1$ inclusive, so each $x$ is in exactly
	one set of the partition.
	Thus, $(s_{n_k})_{k\ge 2}$ is exactly the sequence of partial sums
	of the series
	\[\sum_{n\ge 1} \frac{1}{n^p},\]
	which, for $p > 1$ is known to converge!
	By the earlier comparison argument, this means
	$\sum a_n$ also converges.\\
	
	It remains to prove the converse, that is, when $p\ge 1$, $\sum a_n$
	does not converge. We show that it diverges to $+\infty$
	with a similar comparison test argument.
	This time, for $n\ge 2$, define the sequence
	\[c_n = \frac{1}{2^r\cdot r^p}\]
	where $r$ is the integer such that $2^{r-1} < n \le 2^r$.
	We see that $\log_2 n \le r$, so
	\[a_n = \frac 1n\cdot\frac{1}{(\log_2n)^p} \ge
	\frac{1}{2^r\cdot r^p} = c_n.\]
	Since $c_n$ is always nonnegative and $a_n\ge c_n$ for $n\ge 2$,
	the comparison test says that $\sum_{n\ge 2}a_n = +\infty$
	follows from $\sum_{n\ge 2}c_n = +\infty$.
	
	Again, $c_n$ being nonnegative over all $n$ means that the partial sums
	$(t_n)$ of $(c_n)$ have a limit $\ell$, so any subsequence of $(t_n)$
	also has limit $\ell$. To show that $\ell = +\infty$,
	we consider the subsequence $(t_{m_k})_{k\ge2}$ given by $m_k = 2^{k}$:
	\begin{align*}
		s_{2^k} &= \sum_{j=2}^{2^k} c_j =
		\sum_{r=1}^{k} \sum_{j=2^{r-1}+1}^{2^r} c_j\\
		&= \sum_{r=1}^{k} \sum_{j=2^{r-1}+1}^{2^r} \frac{1}{2^r\cdot r^p}
		= \sum_{r=1}^k \frac{1}{2^r\cdot r^p}\cdot \frac{1}{2^{r-1}}\\
		&= \sum_{r=1}^k \frac{1}{2r^p}.
	\end{align*}
	This subsequence is thus half the partial sums of the
	series $\sum \frac1{n^p}$. Since, for $p\le 1$, $\sum \frac1{n^p}$
	diverges to $+\infty$, the scaled version $\sum\frac1{2r^p}$
	also diverges to $+\infty$.
	Thus, $\sum_{n\ge2} c_n$ also diverges to $+\infty$, and at this point
	the conditions of the comparison test are fulfilled,
	so $\sum_{n\ge2} a_n = +\infty$ is not a convergent series.
\end{proof}

\begin{remark}
	This solution relies on knowing which values of $p$
	make $\sum \frac 1{n^p}$ converge, which we proved using integration;
	however, we can prove this without integration as well
	by a similar argument with bounding via powers of two
	(analogous to our proof of divergence of $\sum \frac 1n$).
\end{remark}

\begin{proof}[Solution 2]
	Just to show I know how the integral test works.
	Since $f(x) = \frac{1}{x(\log x)^p}$ is decreasing as $x$ increases
	(because $x$ and $(\log x)^p$ are both increasing in $x$),
	we can get the bounds
	\[\int_{2}^N f(x)\,dx \le
	\sum_{n=2}^N f(n) \le f(2) + \int_{3}^N f(x-1)\,dx
	= f(2) + \int_2^N f(x)\,dx\]
	for any integer $N > 2$.
	If $\int_2^N f(x)\,dx$ diverges as $N$ goes to infinity,
	then $\sum_{n\ge 2} f(n)$, being bounded below by $\int_2^\infty f(x)\,dx$,
	must also diverge.
	Similarly, if the integral converges, then
	$\sum_{n\ge 2} f(n)$ is bounded above by $f(2) + \int_2^\infty f(x)\,dx$,
	so the partial sums are increasing (because $f(n) > 0$ for $n \ge 2$)
	and bounded, hence the series converges.
	In other words, the convergence of the series
	is equivalent to the convergence of $\int_2^\infty f(x)\,dx$.
	
	Now, we compute the integral to determine when it converges.
	Let $u = \log x$ so that $du = \frac{dx}x$:
	\begin{align*}
		\int_2^\infty \frac{1}{x(\log x)^p}\,dx &=
		\lim_{N\to\infty}\int_2^N \frac{1}{x(\log x)^p}\,dx\\
		&= \lim_{N\to\infty}\int_{\log 2}^{\log N} \frac{1}{u^p}\,du\\
		&= \lim_{N\to\infty} \begin{cases}
			\frac{(\log N)^{1 - p}}{1 - p} - \frac{(\log 2)^{1 - p}}{1 - p}
			& \text{if } p \ne 1\\
			\log \log N - \log \log 2 & \text{if } p = 1
		\end{cases}
	\end{align*}
	Since $\log N\to\infty$ as $N\to\infty$,
	if $1 - p > 0$, $(\log N)^{1-p}$ grows arbitrarily large
	(so the integral diverges),
	while if $1 - p < 0$, then $(\log N)^{1-p}$ grows arbitrarily small
	(so the integral converges).
	Furthermore, if $p = 1$, then $\log\log N$ also goes to infinity.
	Thus, the integral converges if and only if $p > 1$,
	which means the series also converges if and only if $p > 1$.
\end{proof}


\begin{problem}{15.6a}
	Give an example of a divergent series $\sum a_n$
	for which $\sum a_n^2$ converges.
\end{problem}

\begin{soln}
	Let $a_n = \frac 1n$ for all natural $n$.
	Then we have already shown that $\sum \frac 1n$ diverges to $+\infty$,
	while $\sum \frac 1{n^2}$ converges.
	(In fact $a_n = \frac{1}{n^p}$ for any $p\in(\half, 1]$
	easily does the trick).
\end{soln}

\begin{problem}{15.7a}
	Prove that if $(a_n)$ is a decreasing sequence of real numbers
	and if $\sum a_n$ converges, then $\lim na_n = 0$.
\end{problem}

\begin{soln}
	Pick any real number $\eps > 0$; we will show that
	$|na_n - 0| < \eps$ is eventually true.
	To start with, we note that $\sum a_n$ converging means that it satisfies
	the Cauchy criterion, i.e.
	\[\text{there exists an $N$ such that } n > m > N\implies
	\abs{\sum_{k=m}^n a_k} < \frac{\eps}{2}.\]
	
	\begin{claim*}
		All terms of the sequence $(a_n)$ are nonnegative.
	\end{claim*}
	\begin{proof}
		Because $\sum a_n$ converges, $\lim_{n\to\infty} a_n = 0$.
		Since $a_n$ is a decreasing sequence,
		it converges to $\inf\{a_n\}$, which is thus zero;
		this implies that $a_n \ge 0$ for all $n$,
		as $\inf\{a_n\} = 0$ is, by definition,
		less than or equal to each of the $a_n$s.
	\end{proof}
	
	Now, since $a_k \le a_n$ whenever $k \ge n$, we can say that
	\[(n - N)a_n = \sum_{k=N+1}^n a_n \le \sum_{k=N+1}^n a_k =
	\abs{\sum_{k=N+1}^n a_k} < \frac{\eps}{2}\]
	whenever $n > N$. To be clear, the left inequality
	is from $(a_n)$ being decreasing, the equality of sums is
	from $(a_n)$ always being nonnegative,
	and the right inequality comes from the Cauchy criterion
	because $n > N+1 > N$.
	
	Since $N$ is just a constant, we can say that
	\[\lim_{n\to\infty} a_n = \lim_{n\to\infty} Na_n = 0,\]
	so eventually $Na_n < \frac\eps2$ as well, say when $n > M$.
	This means that
	\[n > \max(M, N)\implies
	|na_n - 0| = na_n = (n - N)a_n + Na_n < \frac\eps2 + \frac\eps2 = \eps.\]
	Thus, $|na_n - 0|$ is eventually smaller than any positive real number,
	so $\lim_{n\to\infty} na_n = 0$.
\end{soln}

\pagebreak
\begin{problem}{17.2}
	Let $f(x) = 4$ for $x \ge 0$, $f(x) = 0$ for $x < 0$,
	and $g(x) = x^2$ for all $x$. Thus $\operatorname{dom}(f)
	= \operatorname{dom}(g) = \RR$.
	\begin{enumerate}[\emph{\textbf{(\alph*)}}]
		\ii Determine the following functions:
		$f + g$, $fg$, $f\circ g$, $g\circ f$.
		Be sure to specify their domains.
		\ii Which of the functions
		$f$, $g$, $f + g$, $fg$, $f\circ g$, $g\circ f$ is continuous?
	\end{enumerate}
\end{problem}

\begin{soln}
	For (a):
	\begin{itemize}
		\ii $(f + g)(x) = x^2$ for $x < 0$ and $4 + x^2$ for $x \ge 0$.
		\ii $(fg)(x) = 0\cdot x^2 = 0$ for $x < 0$ and $4x^2$ for $x\ge 0$.
		\ii $(f\circ g)(x) = f(0) = 0$ for $x < 0$ and $f(4) = 16$ for $x\ge 0$
		\ii $(g\circ f)(x) = g(x^2) = 4$ for all $x$ because
		$x^2 \ge 0$ for all real $x$.
	\end{itemize}
	
	For (b), I claim that $g$, $fg$, and $g\circ f$ are continuous
	while the rest of the functions aren't.
	First, to see that these are continuous:
	\begin{itemize}
		\ii $g(x)$ is continuous because $\lim_{n\to\infty} x_n = x$ implies
		\[\lim_{n\to\infty} x_n^2 =
		\lim_{n\to\infty} x_n \cdot \lim_{n\to\infty} x_n = x\cdot x = x^2.\]
		\ii We show that $(fg)(x)$ is continuous at $r$
		by considering three cases based on the sign of $r$.
		If $r > 0$, then any sequence $(x_n)$ converging to $r$ is eventually
		positive (since $|x_n - r| < |r|$ eventually).
		This means $f(x_n) = 4x_n^2$ is eventually true, so
		\[\lim_{n\to\infty} (fg)(x_n) = \lim_{n\to\infty} 4x_n^2 = 4r^2
		= (fg)(r).\]
		Similarly, if $r < 0$, then any sequence $(x_n)$ converging to $r$
		is eventually negative and thus $(fg)(x_n)$ is eventually always zero;
		thus in this case
		\[\lim_{n\to\infty} (fg)(x_n) = \lim_{n\to\infty} 0 = 0 = (fg)(r).\]
		Finally, if $r = 0$, then observe that, for all $\eps > 0$,
		\[|x - r| = |x - 0| < \frac{\sqrt\eps}{2} \implies
		|(fg)(x) - (fg)(r)| = |(fg)(x)| < \eps\]
		because, if $x \ge 0$, then $|(fg)(x)| = 4x^2
		< 4(\frac{\sqrt\eps}{2})^2 = \eps$,
		and, if $x < 0$, then $|(fg)(x)| = 0 < \eps$,
		so $fg$ is continuous at $0$ by the $\eps$-$\delta$ definition
		of continuity (with $\delta = \sqrt\eps$ here).
		\ii Finally, $g\circ f$ is constantly $4$, so it is clearly continuous
		at any real $r$, since, for any sequence $(x_n)$ with limit $r$,
		\[\lim_{n\to\infty} (g\circ f)(x_n) = \lim_{n\to\infty} 4 =
		4 = (g\circ f)(r).\]
	\end{itemize}
	
	Now, we show that the other three functions are all discontinuous at $0$.
	In each case, we need to find a sequence of real numbers $(x_n)$
	that converges to zero such that
	the sequence $(f(x_n))$ doesn't converge to $f(0)$.
	We show that taking $x_n = -\frac1n$ suffices;
	the important property is that $x_n < 0$ for all $n$.
	\begin{itemize}
		\ii For $f$, observe that $f(x_n) = 0$ because $x_n < 0$ for all $n$, so
		\[\lim_{n\to\infty} f(x_n) = 0 \ne 4 = f(0).\]
		\ii For $f+g$, $(f+g)(x_n) = x_n^2$ because $x_n < 0$ for all $n$, so
		\[\lim_{n\to\infty} (f+g)(x_n) = (\lim_{n\to\infty} x_n)^2 = 0
		\ne 4 = f(0).\]
		\ii For $f\circ g$, $(f\circ g)(x_n) = 0$ because $x_n < 0$ for all $n$,
		so
		\[\lim_{n\to\infty}(f\circ g)(x_n) = 0 \ne 16 = f(0).\]
	\end{itemize}
	This is enough to prove that these three functions are all
	not continuous at zero, hence not continuous.
\end{soln}

\begin{problem}{17.4}
	Prove the function $\sqrt x$ is continuous on its domain $[0,\infty)$.
\end{problem}

\begin{soln}
	To prove that $f(x) = \sqrt x$ is continuous at any $r\in[0,\infty)$,
	we want to prove that, for each $\eps > 0$, there exists $\delta$ such that
	\[|x - r| < \delta \implies |\sqrt x - \sqrt r| = |f(x) - f(r)| < \eps.\]
	
	If $r > 0$, then, taking $\delta = \eps\sqrt r$,
	we can see by difference of squares that
	\[\abs{\sqrt x - \sqrt r} = \frac{\abs{x - r}}{\sqrt x + \sqrt r}
	< \frac{\delta}{\sqrt x + \sqrt r} < \frac{\delta}{\sqrt r} = \eps.\]
	
	If $r = 0$, then take $\delta = \eps^2$.
	Note that $f(x)$ is increasing because, if $a = u^2$ and $b = v^2$
	where $u = \sqrt a = f(a)$ and $v = \sqrt b = f(b)$ are real numbers,
	\[a \le b \iff u^2 - v^2 = (u + v)(u - v) \le 0,\]
	which necessitates $u - v \le 0 \iff u \le v$ because $u + v \ge 0$.
	Thus, we have
	\[x = |x - 0| < \delta = \eps^2 \implies
	\sqrt x = |\sqrt x - 0| = |f(x) - f(0)| < \sqrt\delta = \eps,\]
	as desired.
\end{soln}

\begin{problem}{17.9c}
	Prove that $f(x)$, defined as $f(r) = r\sin(\frac 1r)$ for $r\ne 0$
	and $f(0) = 0$, is continuous at $0$.
\end{problem}

\begin{soln}
	For any $\eps > 0$, let $\delta = \eps$.
	Then, to prove that $f$ is continuous at zero,
	we need to show that, for every $\eps = 0$,
	\[|x| = |x - 0| < \delta = \eps \implies |f(x)| = |f(x) - f(0)| < \eps.\]
	Indeed, if $|x| < \delta$, then
	\[\abs{f(x)} = \abs{x\sin\left(\frac 1x\right)} =
	\abs{x}\cdot\abs{\sin\left(\frac 1x\right)} \le
	\abs x\cdot 1 < \delta = \eps,\]
	so the desired implication holds and $f(x)$ is continuous at $0$.
\end{soln}

\pagebreak
\begin{problem}{17.13a}
	Let $f(x) = 1$ when $x$ is rational and $f(x) = 0$ when $x$ is irrational.
	Show $f$ is discontinuous at every $x$ in $\RR$.
\end{problem}

\begin{soln}
	Let $r\in\RR$ be any real number. To show that $f$ is discontinuous at $r$,
	we need to show that there exists an $\eps > 0$ such that,
	regardless of the value of $\delta > 0$, the implication
	\[|x - r| < \delta \implies |f(x) - f(r)| < \eps\]
	does not hold.
	
	Pick $\eps = 1$. Then, we know that, for every pair of real numbers
	$a$ and $b$ with $a < b$, there exists a rational number in $(a, b)$
	and an irrational number in $(a, b)$
	(the former is just a result of denseness of $\QQ$ in $\RR$,
	and the latter was proven in a homework).
	In this case, taking $a = r - \delta$ and $b = r + \delta$,
	we see that there is a rational number $q\in(r - \delta, r + \delta)$
	and an irrational number $y\in(r - \delta, r + \delta)$.
	By definition, $|q - r| < \delta$ and $|y - r| < \delta$,
	and $f(q) = 1$ and $f(y) = 0$. Thus, either $f(r) \ne f(q)$
	or $f(r) \ne f(y)$. If $f(r)\ne f(q) = 1$, then $f(r) = 0$, so we get
	\[|f(r) - f(q)| = |0 - 1| = 1 \ge 1 = \eps,\]
	and if $f(r) \ne f(y) = 0$, then $f(r) = 1$,
	so similarly $|f(r) - f(y)| = 1 \ge \eps$.
	Thus, we have shown that the implication
	\[|x - r| < \delta \implies |f(x) - f(r)| < \eps\]
	fails for $\eps = 1$, which is enough to prove that
	$f$ is discontinuous at any real number $r$.
\end{soln}

\end{document}